\part*{Part VII --- Learning, and the fractal}
\addcontentsline{toc}{part}{Part VII --- Learning, and the fractal}

\bigskip
\noindent
\emph{This part cashes the title. A learner holding a graded formula has two
degrees of freedom, not one: it can change the formula, or it can change the
value. The two live in different spaces, obey different geometries and need
different search methods, and the resulting structure is a fibration rather than a
space. The part then follows the construction upward --- learners compose, and the
composite is itself a learner --- and closes on what an observation discipline
implies about what a mortal thing can know. Sources: \cite{ScientistNote} for the
lattice and the ecology, \cite{ComposeNote} for composition,
\cite{EngineNote} for individuation, \cite{ScopeNote} for the counting.}
\bigskip

\section{The hypothesis space is a fibration}
\label{sec:fibration}

Recall the two structures now in play.

The formulae of a generated logic form the relaxation lattice, ordered by logical
strength and metrised by modal depth: formulae agreeing to depth $n$ are within
$2^{-n}$. This is an ultrametric, so the balls are nested or disjoint, every point
is a centre, and there is no betweenness. There is no gradient
(\S\ref{sec:hypothesis}). Search is backtracking, priced by how far one had to
back up, and Confinement says small steps never leave the initial ball.

The values of a graded clause form the algebra $\V$, which for the learning
disciplines of Part VI is $[0,1]$ or $\R_{\ge0}$ --- ordered, continuous, and
navigable by exactly the local methods the ultrametric denies.

\begin{observation}[Two degrees of freedom]
\label{obs:twodof}
A graded hypothesis is a pair: a formula and a value. Revising the value is not a
move in the ultrametric at all --- it is a move of distance zero, since the formula
is unchanged. So the hypothesis space is not a metric space with extra decoration;
it is a fibration, with the ultrametric tree of formulae as base and a
$\V$-valued fibre over each point.
\end{observation}

\begin{figure}[h]
\centering
\begin{tikzpicture}[font=\small]
% fibres, drawn first so the tree sits on top
\foreach \x in {-2.7,-0.9,0.9,2.7} { \draw[gray!60] (\x,0) -- (\x,2.5); }
% base tree, hanging below the leaf line
\foreach \x in {-2.7,-0.9,0.9,2.7} { \fill (\x,0) circle (2.2pt); }
\draw[thick] (-2.7,0) -- (-1.8,-0.9) -- (-0.9,0);
\draw[thick] (2.7,0) -- (1.8,-0.9) -- (0.9,0);
\draw[thick] (-1.8,-0.9) -- (0,-1.8) -- (1.8,-0.9);
\node[font=\footnotesize,align=center] at (0,-2.5)
  {base: the ultrametric tree of formulae\\ backtracking search; a move discards evidence};
% the two kinds of move
\draw[very thick,-{Stealth}] (0.9,0.35) -- (0.9,2.1);
\node[font=\footnotesize,anchor=west] at (1.05,1.25) {fibre move};
\draw[very thick,-{Stealth}] (-2.62,0.85) -- (-0.98,0.85);
\node[font=\footnotesize,anchor=south] at (-1.8,0.92) {base move};
\node[font=\footnotesize,anchor=west,align=left] at (3.15,1.5)
  {fibres: $\V$-values\\ gradient-like search;\\ a move accumulates\\ evidence};
\end{tikzpicture}
\caption{The hypothesis space of a graded learner. Confinement is a theorem about
the base; a fibre move has base-distance zero, so it cannot leave a ball --- but it
can move anywhere inside one.}
\label{fig:fibration}
\end{figure}

This is the correct statement of the claim that graded where-clauses are a
learning instrument, and it is more modest and more useful than the version one is
tempted to assert.

\paragraph{What grading does not do.} It does not escape Confinement. Confinement
is a theorem about the base, and a fibre move has base-distance zero, so a learner
that only ever revises values is confined \emph{absolutely} --- it never leaves the
ball it started in, because it never leaves the point it started at. The forager
of \S\ref{sec:d3} is a case in point: it learns to stop opening, which saves its
life, but it never acquires a better test.

\paragraph{What grading does do.} It makes the interior of each ball navigable.
Before grading, the only way to change behaviour was to change the formula --- a
discrete jump, priced by backtracking depth, which invalidates every confirmation
obtained below that depth. After grading, a learner can change its behaviour
continuously while holding its formula, and its evidence, fixed. In the harsh
world of \S\ref{sec:d3} that difference is the difference between $.0029$ and
survival.

\begin{principle}[Two searches, one learner]
\label{prin:twosearch}
The base requires backtracking search; the fibre admits gradient-like search. A
learner needs both, and the framework prices them differently: a base move
discards evidence, and a fibre move accumulates it.
\end{principle}

There is a pleasing symmetry in the failure modes. A learner that only moves in
the fibre is confined to one theory and can only tune it --- it is the forager, safe
and stuck. A learner that only moves in the base is the ladder-climber of
\S\ref{sec:d2}, which can acquire deeper theories but pays for each and, as
measured, may be made \emph{worse} by depth if it holds its other commitments
fixed. Both pathologies are visible in the numbers, and neither is visible without
the other coordinate.

\begin{remark}[Where the two searches meet]
The graded fixed point of \S\ref{sec:engines} is the point of contact. $\mu$ is
search and $\nu$ is survival, and both are fibre-side operations whose availability
depends on the value algebra --- free in a quantale, a convergence condition over
$\R_{\ge0}$, a budget over $\Cx$. So how far a learner can search \emph{within} a
theory is settled by the same choice that settles whether it can state a
refutation. Part IV's negation criterion and Part V's fixed-point criterion are the
same criterion.
\end{remark}

\section{What escapes the base}
\label{sec:escape}

If small moves cannot leave a ball, something else must. The corpus identifies
three operators that can, and all three are structural rather than epistemic.

\paragraph{Recombination.} Reproduction by crossover on quoted code is a base move
of arbitrary size. Because the germ is a quotation --- inert, transmissible,
inspectable --- and because the normal form of a process is a tree of communicating
pairs, crossover can swap material at homologous loci, producing a formula at a
distance no individual revision could reach. The unit of inheritance is a
\emph{redex}: a gene is an interaction, not a structure, and linkage is derived
rather than posited. And crucially, a revision move and a crossover are \emph{one
move} carried along the characteristic-formula map, applied at two timescales and
distinguished by accessible distance \cite{ScientistNote}.

\paragraph{Composition.} Two learners in parallel constitute a third. This is the
subject of \S\ref{sec:compose}, and it is the first of the escapes that requires no
new syntax at all --- parallel composition is already in the calculus.

\paragraph{Death.} A population whose members are confined to different balls
explores what no member can. This is why the framework insists that a population,
not an individual, is the unit of continual learning: not because populations are
better, but because Confinement makes a single learner structurally incapable of
leaving its basin, and because mortality supplies the retirement criterion that a
purely epistemic account has to invent.

\section{Learners compose, and the composite is a learner}
\label{sec:compose}

A learner, on this account, is a distribution over a population of processes
together with resource reservoirs; the population constitutes a \emph{namespace};
and learners compose by parallel composition \cite{ComposeNote}.

\begin{theorem}[Factorisation]
\label{thm:factor}
Namespace disjointness, absence of a spanning redex, and factorisation of the
composite's distribution as a product are equivalent. Graded separation at $k=0$
\emph{is} probabilistic independence, and the defect of the comparison map is
exactly the interaction.
\end{theorem}

Overlap is therefore the dial, and it is a vector: one grade per typed namespace
--- metabolic, source, problem, mating. The familiar ecological relations
(neutralism, competition, predation, mutualism, commensalism, theft) fall out as
sign conditions on that vector. And one asymmetry is forced:

\begin{proposition}[Cooperation is epistemic, competition is metabolic]
\label{prop:coopepis}
Overlap confined to the metabolic and source namespaces has non-positive total
gain: there is no metabolic mutualism. Positive sum lives only in the problem
namespace, via amortised perception --- which is Principle~\ref{prin:readdrive}
again, one level up, in the form of the question of who is billed.
\end{proposition}

Two further results are worth carrying, because they bear on the design of
multi-agent systems rather than on biology.

\paragraph{Non-interference is unstable, and modularity is purchasable.}
Independently manufactured names can coincide, so disjointness is luck. Rooting
each learner's namespace at its own quotation makes disjointness a theorem, via
freshness-by-quotation. The price is a naming discipline.

\paragraph{Competitive dilution is regressive.} At a shared rate-limited source,
with shares proportional to inverse cost, each learner retains exactly
$c_{\mathrm{other}} / (c_{\mathrm{self}} + c_{\mathrm{other}})$ of the deep
inquiry it could afford alone --- independent of the reservoir, the rate, and the
price. In the worked composite, the learner with the incomplete theory loses
$66.9\%$ of its affordable depth while the one with the complete theory loses none
of its depth and only a third of its assays; at exhaustion the latter holds
$5452$ against the former's $1397$. Competitive exclusion in epistemic currency,
requiring no contact at all.

\section{Testimony, and a counterfeiting attack}
\label{sec:testimony}

Once values are graded, learners can exchange \emph{values} and not only
observations. In a probabilistic logic, revision adds evidence counts, so
confidence is superadditive; and Proposition~\ref{prop:coopepis} makes metabolism
subadditive. The programme's slogan therefore becomes an inequality rather than a
sentiment.

\begin{observation}[Belief inflation is the epistemic analogue of minting]
\label{obs:minting}
Revision assumes independent evidence. Two learners who both heard from a third
double-count, and nothing in the framework forbids it. Tokens are conserved;
confidence is not. So there is a counterfeiting attack in the problem namespace
with no analogue in the metabolic one, and the plagiarism equilibrium becomes a
condition on the grading vector rather than a worry.
\end{observation}

This is stated and not developed. It requires carrying strength and confidence as
a pair rather than collapsing them to a product, which needs a lax algebra with a
canonical aggregation order, and that is a paper of its own.

\begin{remark}[A belief is a message]
In the calculus the belief sits on a channel as data. Testimony and theft of a
hypothesis are therefore the same operation as predation, differing only in which
namespace the name lives in. That is either an unfortunate consequence or an
accurate model of intellectual life.
\end{remark}

\section{The fractal}
\label{sec:fractal}

The construction repeats. An individual is a learner; a population of learners is
a learner; a lineage is a learner over a longer timescale; a composite of
populations is a learner again. What makes this more than a rhetorical spiral is
that the generating step is identified and priced.

\paragraph{The generating step is the medium.} Communication in the calculus is
perfect, without loss of generality, because error is modelled by interposition: a
communication becomes a send, a channel process, and a receive, where the
interposed process may forward, drop, corrupt, read or convert. Those interposed
processes form a namespace of their own --- and \emph{the channel namespace at one
level is the endpoint namespace at the next}. Well-founded downward by the
no-self-code theorem, bounded upward by the profitable radius. The outside of one
level is the inside of the next \cite{EngineNote}.

\begin{figure}[h]
\centering
\begin{tikzpicture}[font=\footnotesize,
  lv/.style={draw,rounded corners=2pt,inner sep=6pt}]
\node[lv,minimum width=1.6cm,minimum height=1.0cm] (a) at (0,0) {};
\node at (0,0) {endpoints};
\node[lv,minimum width=3.4cm,minimum height=2.0cm] (b) at (0,0) {};
\node at (0,0.75) {media $\mathsf{Chn}$};
\node[lv,minimum width=5.2cm,minimum height=3.0cm] (c) at (0,0) {};
\node at (0,1.25) {endpoints, one level up};
\node[lv,minimum width=7.0cm,minimum height=4.0cm] (d) at (0,0) {};
\node at (0,1.75) {media, one level up};
\draw[-{Stealth}] (3.8,-1.6) .. controls (5.2,-1.0) and (5.2,1.0) .. (3.8,1.6);
\node[anchor=west,align=left] at (5.0,0)
  {the outside of one level\\ is the inside of the next};
\end{tikzpicture}
\caption{The generating step. Well-founded downward by the no-self-code theorem;
bounded upward by the profitable radius.}
\label{fig:tower}
\end{figure}

\paragraph{Individuation is error correction.} A region may close its namespace ---
become an individual --- exactly when its boundary is rich enough relative to its
radius: $h(B) \ge (\epsilon\theta/\rho)\,r$, with $\epsilon$ the per-hop error
rate, $\rho$ the repair throughput and $\theta$ the toll. Polynomial growth
recovers a critical radius $r^{*} = \sqrt{\rho/\epsilon\theta}$; an expander makes
it linear. Nature rides the boundary between centralisation and decentralisation
because centralisation compresses and decentralises poorly under error, and error
correction is the principle that says when to close over a namespace.

\paragraph{Individuals are porous, and porosity is not optional.} Perfect closure
is death: no harvest, no assay. An individual has a closed interior namespace and
an open boundary namespace, and the boundary channels must have redex pathways to
the interior or they are stranding sites. Deep integration is exponentially
fragile in the number of hops, which is why ownership is what buys depth.
Individuation is a fixed point, and the circularity is the fractality.

\paragraph{What this costs in copies.} Following \S\ref{sec:d6}: a multi-skilled
individual is a composite of ecologies-as-learners, each skill an ecology in its
own right, and the copy number is a series indexed by stratum rather than an
integer. Depth is exponentially fragile and width is linear, so buying depth
obliges you to buy copies --- the two axes are not orthogonal, which corrects an
earlier statement in this corpus.

\begin{observation}[The loop closes]
\label{obs:loop}
The individual at the top of one tower is the scientist at the bottom of the next.
There is no distinguished level, and no level at which the observation discipline
stops mattering: each stratum has its own namespace, its own budget, and its own
logic generated from its own term structure.
\end{observation}

\section{An aggregate mind assembles several disciplines}
\label{sec:aggregate}

\noindent\emph{In plain terms: a mind does not have one logic. It has several,
running side by side, seeing different things --- and this is not an accident of
biology but something the framework forces.}

Everything in \S\ref{sec:compose} composed learners that shared an observation
discipline and differed only in what they knew. Nothing requires that. A learner
is a distribution over a population, and different sub-populations may carry
different dial settings. Parallel composition does not care.

\subsection{Two games}

Consider the card game \emph{Set}. Each card carries four attributes with three
values each, and a legal set is three cards on which every attribute is either
all-the-same or all-different. What does it take to see one? It takes counting
occurrences of abstract categories: three reds, three distinct shapes, three
distinct textures. And counting is exactly what a fully structural discipline
cannot do.

\begin{observation}[Contraction is what makes counting impossible]
\label{obs:counting}
Contraction corresponds to idempotence of the container
(Principle~\ref{prin:dial}). A learner whose witnesses are collected in a set
silently discards multiplicity: it can see \emph{that there are red ones}, and it
cannot see \emph{that there are three}. Seeing three requires a container that
withholds contraction --- a bag --- and the logic that container generates is the
multiplicative fragment of \S\ref{sec:linear}. \emph{Set} is a game one can only
play in the linear discipline.
\end{observation}

Now consider poker. A player needs the same counting --- outs, pot odds, how many
of a rank remain --- and needs something else entirely: to read the table. Who is
playing with whom, which subset of players could between them take the pot without
the rest, whose betting is advocacy and whose is noise. That is not a question
about multiplicity. It is a question about which coalitions are actionable, and
the discipline that generates it is the semitopological one of
\S\ref{sec:semitop}.

A poker player is therefore a learner that composes two observation disciplines
with \emph{different containers}. And here the framework says something sharper
than ``that is convenient''.

\begin{observation}[The disciplines cannot be merged, only composed]
\label{obs:nomerge}
There is no single container that is both. Merging would require a distributive
law between the two collection monads, and the no-go theorems for distributive
laws turn on exactly the equations that separate them --- commutativity and
idempotence \cite{ZwartMarsden19}, the same two coordinates
Principle~\ref{prin:dial} is built on. So an aggregate mind cannot have one
unified logic that sees everything its parts see. It must run several disciplines
in parallel and let them exchange results through a shared namespace, which is
composition in the sense of \S\ref{sec:compose} and is available with no new
syntax.
\end{observation}

We would rather understate this than overstate it, so: the no-go theorems are
about specific pairs of monads and their proofs are not uniform, and we have not
checked the particular pair at issue. What is safe to say is that merging
containers is exactly the operation those theorems obstruct, that the obstruction
is generic, and that the framework therefore predicts modularity rather than
having to assume it. If that is right, then the differentiated capabilities one
observes in minds --- a faculty for quantity, a faculty for social reading, a
faculty for space --- are not an engineering compromise inherited from evolution.
They are what an aggregate observer has to look like.

\subsection{Capabilities are separately present}

Disciplines are separately generated, separately priced, and separately acquired.
It follows immediately that they can be separately \emph{absent}, and that the
absence has a characteristic signature.

A learner missing a discipline is not globally impaired. It is precisely blind to
the distinctions that discipline generates, and fully competent on the rest. By
Principle~\ref{prin:fewerkinds} it also does not experience the blindness as
blindness: it reports a world with fewer kinds in it than the world has, and the
report is internally consistent. And by the legibility results of
\S\ref{sec:d5}, what such a learner loses in a population it cannot read is
\emph{coverage}, not accuracy --- it declines to answer rather than answering
wrongly, which from the outside looks like reticence and from the inside looks
like a world in which those questions simply do not arise.

\begin{remark}[A speculation, flagged as such]
\label{rem:autism}
It is tempting to read some of the observed variation in human social cognition
this way --- to conjecture that what differs is the availability or the weighting
of something like the coalition discipline, rather than any general capacity, and
that this would explain why such variation so often coexists with unimpaired or
superior performance on exactly the counting-and-structure disciplines the same
framework generates.

Two things should be said immediately. First, the framework admits a second
reading of the same phenomenon, and it is not the deficit reading: two learners
may both have a coalition discipline and disagree about which coalitions count as
actionable --- a different container, not a missing one. That reading predicts
\emph{symmetric} mutual illegibility rather than a one-way deficit, since each
population is reading the other with the wrong opens, and \S\ref{sec:d5}'s
counter-move result says the failure would show up as loss of coverage on both
sides rather than error on either. The empirical literature contains a well-known
argument of exactly that shape, and we note that the framework's natural
prediction sits with it rather than against it.

Second, we are in no position to adjudicate. This is a conjecture stated in the
framework's vocabulary to show that the vocabulary has empirical bite, and it is
not a clinical claim; nothing in this document was built to make claims about
people. We record it because a reader will think of it, and because the two
readings are genuinely distinguishable --- one predicts asymmetric failure, the
other symmetric --- which makes it a question that could be settled rather than
argued.
\end{remark}

\subsection{The composite is a learner again}

Note finally that this is the fractal of \S\ref{sec:fractal} in another costume. A
poker player composing a counting discipline with a social one is a learner whose
sub-learners run different logics, exchanging through the problem namespace ---
which is where, by Proposition~\ref{prop:coopepis}, the positive-sum lives. The
composite is itself a learner, with its own budget and its own confinement, and it
can be composed again.

\begin{figure}[h]
\centering
\begin{tikzpicture}[font=\small,
  d/.style={draw,rounded corners=2pt,align=center,inner sep=5pt,minimum width=3.9cm,minimum height=1.25cm},
  ns/.style={draw,dashed,rounded corners=2pt,align=center,inner sep=5pt,minimum width=3.4cm}]
\node[d] (lin) at (-2.6,1.5) {$\Coll = \Mset$\\[1pt]\footnotesize counting; \emph{Set}};
\node[d] (coa) at (2.6,1.5) {$\Coll = \Open(P)$\\[1pt]\footnotesize reading the room};
\node[ns] (ns) at (0,-0.4) {shared problem namespace};
\draw[-{Stealth}] (lin.south) -- (ns.north west);
\draw[-{Stealth}] (coa.south) -- (ns.north east);
\node[align=center,font=\footnotesize] at (0,-1.5) {the composite: poker};
\draw[dotted,thick] (-4.8,-2.0) rectangle (4.8,2.4);
\end{tikzpicture}
\caption{Two disciplines, two containers, one learner. There is no container that
is both (Observation~\ref{obs:nomerge}), so the composition happens at the level
of learners.}
\label{fig:aggregate}
\end{figure}

\section{What a learner can afford to see}
\label{sec:afford}

The document ends where it must, which is in epistemology rather than engineering.

Every result above is a result about a \emph{mortal} observer: one with a purse,
whose observations are priced by depth, whose refutations are budget-relative, and
whose undetermined verdicts have two sources it cannot tell apart. Such an
observer does not have a logic in the sense of a fixed vocabulary handed down. It
has an observation discipline, which is generated from the machinery it is made
of and paid for out of the same budget as its food.

The consequences compound in one direction.

Its senses are the structural predicates it can afford, so it perceives shapes to
a certain depth and no further. Its lookahead is the modal nesting it can afford,
which by \S\ref{sec:d2} is usually a poor purchase. Its hypotheses are formulae in
the relaxation lattice, so its theories are confined to the ball it began in
unless recombination or composition move it. Its counting is graded by its
resolution, so it reports fewer kinds than the world contains --- and, by
Principle~\ref{prin:fewerkinds}, does not see that it is doing so. Its refutations
run out when its purse does. And the very distinctions available to it are fixed
by dials it did not choose and cannot inspect from inside.

That is Kant's picture with a ledger attached. The categories are not transcendent
and they are not arbitrary: they are what a mortal, interacting learner can afford,
generated from what it is made of. The framework's contribution is not the
observation --- philosophers have had it for two centuries --- but the arithmetic:
the budget constraint can be \emph{measured}, and the difference between what a
learner sees and what is there can be computed.

\begin{remark}[On the effectiveness of reason]
There is a well-known puzzle about why mathematics should be so effective on the
natural world, and on this account the puzzle is asked across a gap that is not
there. The hypothesis language is generated from the substrate, and adequacy ---
where it holds --- makes it separate exactly what interaction separates. So a
formal language is effective on exactly the domain to which interaction reaches,
because it is the shadow that the access relation casts. Two consequences follow
that are not merely restatements: transfer to non-ancestral domains should be
expected, because adequacy is for interaction \emph{as such} rather than for any
particular interaction; and there should be a predicted boundary, at any relation
that is not interaction across a cut. Both are stated in \cite{ScientistNote} and
neither is established here.
\end{remark}
